% ============================================================
%  ch3_related_work.tex
%  Chapter 3: Related Work
%  Master's Thesis -- SQL MATCH_RECOGNIZE on Pandas DataFrames
% ============================================================
% !TEX root = Thesis.tex
\chapter{Related Work}
\label{chap:related}

This chapter positions the thesis within existing work on SQL
row-pattern recognition and Python-based tabular analytics.
Section~\ref{sec:related-native} reviews systems that provide native
\texttt{MATCH\_RECOGNIZE} or related row-pattern support.
Section~\ref{sec:related-pandas-gap} discusses the gap in Pandas and
similar DataFrame libraries.  Section~\ref{sec:related-contribution}
summarizes how this thesis addresses that gap.

Row pattern recognition was first proposed as a SQL standard feature
by Zemke et al.~\cite{zemke2007pattern} and was officially adopted
in SQL:2016~\cite{ISO2016SQL,iso2016}.
Later standards guidance and technical literature documented its
semantics in more detail~\cite{ISO2021,melton2017sqlrpr,
michels2018new}.  Further work studied practical implementations and
applications~\cite{petkovic2022rpr,petkovic2022fraud}.

The \texttt{MATCH\_RECOGNIZE} clause is supported natively by
several SQL and stream-processing engines, including
Oracle~\cite{OracleMR,oracle2021dwhsg,book,oracle2022},
Snowflake~\cite{snowflake2025match_recognize}, and
Trino~\cite{trino2022,trino2021row}; Apache Flink also exposes row
pattern recognition through Flink SQL~\cite{flink}.
By contrast, PostgreSQL~18 still lists the related SQL features as
unsupported, while DuckDB lists native \texttt{MATCH\_RECOGNIZE} as
planned work~\cite{postgresql18unsupported,duckdb2026roadmap,howe2020matchrecognize,findling2025bringing}.
In systems without native support, users typically emulate
row-pattern matching using window functions, recursive queries,
user-defined aggregates, procedural layers, or query
translation~\cite{TURP2,TURP,cao2012optimizationanalyticwindowfunctions,
lambrecht2025democratize}.
These alternatives can express many of the same ideas, but they are
usually more verbose and harder to maintain than a declarative row
pattern clause.

% ----------------------------------------------------------------
\section{Existing Systems and Approaches}
\label{sec:related-native}
% ----------------------------------------------------------------

Oracle~\cite{OracleMR,oracle2021dwhsg,book},
Trino~\cite{trino2022,trino2021row},
Snowflake~\cite{snowflake2025match_recognize}, and
Flink~\cite{flink} provide SQL-based row-pattern recognition.
Their scope and execution models differ.  Oracle and Snowflake target
database or warehouse workloads.  Trino is a distributed query
engine, while Flink targets streams.  This variety shows the value of
the abstraction, but it does not give Pandas users an in-process
implementation.

The systems reviewed here are representative rather than exhaustive.
They were selected because they provide semantic references,
performance baselines, streaming alternatives, or in-process designs
that are relevant to this thesis.  Other platforms also provide the
clause.  BigQuery documents \texttt{MATCH\_RECOGNIZE}, and Databricks
introduced beta support in Runtime~19.0~\cite{bigquery2025matchrecognize,
databricks2026matchrecognize}.

% ................................
\subsection{Comparative Overview of Existing Implementations}
\label{subsec:related-existing-overview}
% ................................

The SQL standard defines three optional row-pattern
features---R010 (\texttt{FROM}-clause patterns), R020 (window
patterns), and R030 (full aggregate
support)---introduced in Section~\ref{sec:bg-match-recognize}.
In practice, vendor documentation normally
describes the supported syntax instead of declaring conformance by
feature code.  Table~\ref{tab:engine-comparison} therefore describes
the documented scope directly.  It does not claim that all systems
provide identical SQL semantics.

\begin{table}[htbp]
  \centering
  \caption[Scope and deployment of selected row-pattern systems]{Documented scope and deployment of selected row-pattern
    systems and approaches}
  \label{tab:engine-comparison}
  \small
  \setlength{\tabcolsep}{4pt}
  \begin{tabularx}{\textwidth}{l>{\raggedright\arraybackslash}Xll}
    \toprule
    \textbf{System} & \textbf{Documented scope} &
      \textbf{Deployment} & \textbf{License} \\
    \midrule
    Oracle & \texttt{FROM}-clause row patterns &
      On-premises / cloud & Proprietary \\
    Trino & \texttt{FROM} and window row patterns &
      Server / cluster & Apache~2.0 \\
    Snowflake & \texttt{FROM}-clause row patterns &
      Cloud service & Proprietary \\
    Apache Flink & Streaming SQL subset backed by CEP &
      Streaming cluster & Apache~2.0 \\
    DuckDB & No native clause; external research transpiler &
      In-process & MIT \\
    \midrule
    \textbf{This work} & \textbf{Evaluated R010-style subset} &
      \textbf{In-process Python} & \textbf{MIT} \\
    \bottomrule
  \end{tabularx}
\end{table}

% ................................
\subsection{Oracle}
\label{subsec:related-oracle}
% ................................

Oracle introduced SQL row-pattern recognition in Database~12c.  Its
implementation is a main reference for relational
\texttt{MATCH\_RECOGNIZE} semantics~\cite{OracleMR,
oracle2021dwhsg,book}.
The Oracle~21c documentation considered in this thesis places
\texttt{MATCH\_RECOGNIZE} in a table expression; it does not document
the R020 window form.  It also prohibits some nesting and
correlated-subquery forms in \texttt{MEASURES} and \texttt{DEFINE}.
Its documented aggregate support is limited to \texttt{COUNT},
\texttt{SUM}, \texttt{AVG}, \texttt{MAX}, and \texttt{MIN}, without
\texttt{DISTINCT}~\cite{oracle2021dwhsg}.
Oracle uses a proprietary license and requires a separate database
deployment.  This differs from an in-process, open-source package.

% ................................
\subsection{Trino}
\label{subsec:related-trino}
% ................................

Trino is an open-source distributed query engine designed for
parallel, federated query execution across large data
sources~\cite{trino2022,trino2021row}.  Its current documentation
covers \texttt{MATCH\_RECOGNIZE} in the \texttt{FROM} clause and
also documents row pattern recognition in window structures.  This
gives Trino broad documented coverage among the reviewed open-source
systems.  It is still a query server rather than an in-process
DataFrame library.  In this thesis, Trino~473 is a semantic reference
for selected results and a performance baseline.

% ................................
\subsection{Snowflake}
\label{subsec:related-snowflake}
% ................................

Snowflake is a managed cloud data platform.  It documents
\texttt{MATCH\_RECOGNIZE} support, with rich \texttt{DEFINE} and
\texttt{MEASURES} expressions for row-pattern
queries~\cite{snowflake2025match_recognize,hoffa2021snowflake}.
Its documentation also notes implementation-specific behavior.  The
engine uses backtracking, and some pattern and data combinations can
take a long time to execute.  Some row-pattern functions are also
restricted to particular clauses.  Snowflake therefore recommends
selective conditions and bounded quantifiers when possible.  Because
it is a managed cloud warehouse, computational cost also depends on
warehouse size and query workload~\cite{snowflake2025match_recognize}.

% ................................
\subsection{Apache Flink}
\label{subsec:related-flink}
% ................................

Apache Flink is designed for real-time stream processing.  It exposes
\texttt{MATCH\_RECOGNIZE} through Flink SQL, backed by its Complex
Event Processing (CEP) runtime
library~\cite{flink,vsprem2024building,georgiou2025streaming}.
Its documented implementation is a SQL subset oriented toward
streaming use cases.  The current Flink documentation lists
unsupported features including pattern groups, subsequence
alternation, \texttt{PERMUTE}, anchors, exclusions,
\texttt{ALL ROWS PER MATCH}, \texttt{SUBSET}, physical offsets
(\texttt{PREV}/\texttt{NEXT}), and \texttt{DISTINCT}
aggregations~\cite{flink}.  Flink fits continuous event streams, but it
assumes a streaming runtime rather than a lightweight, batch-oriented
DataFrame workflow.

When pattern output is later used in Python-based machine learning,
server and cluster systems may require additional integration between
SQL execution and modeling tools~\cite{
ML_Workflow,Makrynioti2019sql4mlAD}.  In many practical pipelines,
this can introduce data movement, duplicated intermediate data, or
more complicated orchestration steps~\cite{naphade2025evolution}.

% ................................
\subsection{Bringing \texttt{MATCH\_RECOGNIZE} to DuckDB}
\label{subsec:related-duckdb}
% ................................

DuckDB is an in-process analytical database~\cite{raasveldt2019duckdb}.
It needs no separate server, and its Python API can query a Pandas
DataFrame directly from SQL~\cite{duckdb2026pandas}.

At the time of review, \texttt{MATCH\_RECOGNIZE} has not yet been
implemented in DuckDB itself; its roadmap lists the clause as planned
work~\cite{duckdb2026roadmap}.  Lambrecht et al. instead translate
row-pattern queries into SQL that DuckDB already supports.  Their
transpiler rewrites a query into non-recursive and recursive common
table expressions and then executes the translated SQL in
DuckDB~\cite{lambrecht2025democratize,findling2025bringing,mogensen2024compiler}.
This approach expresses pattern variables, ordering, skip behavior,
and measures without adding a native operator to DuckDB.  It still
requires translation and execution through a second engine.

% ................................
\subsection{Optimization Research}
\label{subsec:related-optimization}
% ................................

Zhu et al. study how row-pattern queries over historical relational
data can be optimized~\cite{Zhu2023rowpattern}.  Their approach finds
pattern variables whose conditions can be checked independently.  It
uses relational joins to remove rows that cannot take part in a match,
then runs the NFA matcher on the smaller input.  Cost estimates for the
joins and the NFA guide this choice.

This thesis addresses a different execution environment.  It compiles
conditions and match plans for in-memory arrays, then uses the general
execution path when an optimization does not apply.  Both approaches
have the same broad goal: reduce unnecessary work before or during the
automaton scan without changing the query result.

% ................................
\subsection{Implementation Restrictions Across Systems}
\label{subsec:related-restrictions}
% ................................

The comparison shows three recurring differences~\cite{
oracle2021dwhsg,trino2022,flink,snowflake2025match_recognize,
petkovic2022rpr}.  First, some systems support only the
\texttt{FROM}-clause form, while others also support window patterns.
Second, function and aggregate coverage differs.  Third, empty
matches, exclusions, reluctant quantifiers, and skip modes are not
uniform across implementations.

Oracle, Trino, Snowflake, and Flink also require a database, cluster,
stream runtime, or cloud service.  That setup is reasonable for
workloads already in those environments, but it adds friction when
the data already fits in memory~\cite{8509400,Evaluation,
in-memory_processing}.

% ----------------------------------------------------------------
\section{Pandas and the Python Ecosystem Gap}
\label{sec:related-pandas-gap}
% ----------------------------------------------------------------

Pandas is a popular open-source Python library for data
analysis~\cite{McKinney2010}.  Easy installation, broad adoption, and
tight integration with the Python data-science stack make it a common
choice for tabular analysis without a separate database.
It provides primitives for sorting, filtering, grouping, rolling
windows, shifting columns, and time-series analysis.

But those are low-level building blocks, not a declarative row-pattern
language.  Expressing a multi-step sequence usually means stitching
together \texttt{groupby}, \texttt{shift}, \texttt{rolling}, custom
loops, or intermediate Boolean columns.  The code becomes verbose and
error-prone as patterns lengthen, need overlap control, or refer to
already-matched rows.

The Pandas~2.3.3 API used in this thesis has no native SQL:2016
\texttt{MATCH\_RECOGNIZE}-style primitive~\cite{pandas233api}.
Equivalent logic must therefore be assembled from operations such as
\texttt{groupby}, \texttt{shift}, \texttt{rolling}, and custom
iteration, or sent to another engine.
Other DataFrame libraries, including Dask, Polars, and Modin, are
outside the Pandas-specific implementation and evaluation scope of
this thesis~\cite{Evaluation}.

% ................................
\subsection{Research Gap in the Pandas Ecosystem}
\label{subsec:related-gap-synthesis}
% ................................

The reviewed systems expose a deployment and workflow gap.  Each fits
when the data already lives in its environment or needs distributed or
stream processing.  None is the direct execution path for a Pandas
DataFrame already held in a notebook or script.

DuckDB reduces the setup cost because it runs in-process.  The reviewed
row-pattern approach still translates the query and executes it inside
DuckDB.  The gap addressed here is more specific: a Pandas-native
interface that keeps the declarative SQL structure without a second
engine or a transpilation step.

That sets the scope of this work.  The engine targets in-process,
offline DataFrame analytics and an evaluated R010-style subset with
selected SQL:2016 constructs.  It is not intended to replace a
distributed database or stream processor.

% ----------------------------------------------------------------
\section{Contributions of This Work}
\label{sec:related-contribution}
% ----------------------------------------------------------------

The three objectives in Section~\ref{sec:intro-objectives} lead to
three connected contributions: the engine and its tested semantics,
its optimized execution design, and the evaluation.

\begin{enumerate}
  \item A \textbf{Pandas-native engine} for the evaluated SQL:2016
    R010-style subset.  It supports the evaluated forms of the main
    clauses, output modes, skip modes, navigation functions, and
    aggregates.  Its pattern
    constructs include quantifiers, alternation, grouping, anchors,
    exclusions, and \texttt{PERMUTE}.  It runs in-process and needs no
    database server or cloud warehouse.  Chapter~\ref{chap:methodology}
    presents its five-stage design.

  \item An \textbf{optimized and guarded execution design} with compiled fast paths,
    vectorized predicate evaluation, compact output construction,
    bounded caching, and partition-scoped data handling.  Conditions
    outside a fast path use the general evaluator.

  \item \textbf{Theoretical and empirical evaluation}.  The
    theoretical analysis gives time and space bounds for ordering,
    matching, output construction, and pattern complexity.  The
    empirical analysis evaluates correctness, SQL feature coverage,
    usability, performance, and scalability.  Performance includes
    execution time, throughput, and memory.  The study uses synthetic
    reference cases, Amazon UK Products~2023 up to 2.22~million rows,
    and Amazon Reviews~2023 up to 227.9~million
    rows~\cite{asaniczka2023amazon,hou2024amazonreviews}.
    Chapter~\ref{chap:evaluation} gives the full method and results.
\end{enumerate}

% ----------------------------------------------------------------
\section{Chapter Summary}
\label{sec:related-summary}
% ----------------------------------------------------------------

Oracle, Trino, Snowflake, and Flink show how row patterns work in
database and streaming environments.  DuckDB research shows a
translation-based, in-process alternative.  Optimization research
shows how early filtering can reduce matching work.

Pandas~2.3.3 still has no direct row-pattern interface.  Users must
write procedural DataFrame logic or call another engine.  This leaves
the focused gap addressed by the thesis: declarative row-pattern
recognition inside an in-memory Pandas workflow.

% ................................................................
% END OF CHAPTER
% ................................................................
